\section{Divisor detection in type A}\label{sec:type_a}

We work over $\C$ and use the convention $X_w=\overline{BwB/B}\subset
\GL_n/B$, where $B$ is upper triangular and $\dim X_w=\ell(w)$. For a
normal variety $X$ and a Weil divisor $D$, let $\NonCart_X(D)$ be its
non-Cartier locus. The following statement concerns every Weil divisor,
including divisors whose local classes have infinite order.

\begin{theorem}[Type A divisor detection]\label{ta:main}
For every $w\in S_n$ and every Weil divisor $D$ on $X_w$,
\begin{equation}\label{ta:eq:locus}
 \NonCart_{X_w}(D)=
 \bigcup_{\substack{S\in\Irr\Sing(X_w)\\
                   D\text{ is not Cartier at }\eta_S}}S.
\end{equation}
Equivalently, for every scheme point $p\in X_w$, restriction induces an
injection
\begin{equation}\label{ta:eq:local_injection}
 \Cl(\OO_{X_w,p})\longrightarrow
 \bigoplus_{\substack{S\in\Irr\Sing(X_w)\\p\in S}}
 \Cl(\OO_{X_w,\eta_S}).
\end{equation}
The same statement holds for Schubert varieties in type A partial flag
varieties.
\end{theorem}

The proof constructs detectors for integer differences of divisor
coefficients. Its geometric ingredients are a filtered determinantal
slice construction and a procedure that replaces a horizontal component
of a first-line hyperplane section by a comparison between vertical
components.

\subsection{Coordinates and a first-difference localization}

Put $C_y=ByB/B$, $C^x=B^-xB/B$, $N_{x,y}=X_y\cap C^x$, and
$A_{x,y}=\C[N_{x,y}]$. For $x\leq y$, this slice is integral, normal and
Cohen--Macaulay, with dimension $\ell(y)-\ell(x)$.
The affine chart $\Omega_x=xC^{\id}$ satisfies
\[
 \Omega_x\cap X_y\simeq C_x\times N_{x,y}
 \simeq\A^{\ell(x)}\times N_{x,y},
\]
compatibly with the Schubert strata; see
\cite[Section 3 and Lemma 3.2]{WY08}. The standard normality and rational-singularity results,
and the Cohen--Macaulay property, pass to the slice through this product
\cite[Theorems 2.1.1, 2.2.3 and Proposition 2.2.5]{Bri05}.

For $y\in S_n$, set $r_y(r,k)=\#\{h\leq k:y(h)\geq r\}$.
Then $x\leq y$ exactly when $r_x(r,k)\leq r_y(r,k)$ for every $r,k$.
If $p<q$ and $x(p)<x(q)$, exchanging positions $p,q$ increases the rank
matrix by one precisely on the rectangle
$p\leq k<q,\; x(p)<r\leq x(q)$.
This is the rank-matrix test used below.

In the standard matrix for $C^x$, column $p$ has pivot $1$ in row $x(p)$.
Its free entries are $Z_{x(q),p}$ for $p<q$ and $x(p)<x(q)$, of weights
$e_{x(q)}-e_{x(p)}$. Distinct free entries have distinct root weights.
The equations of $N_{x,y}$ bound the southwest submatrix ranks by
$r_y(r,k)$.

\begin{lemma}\label{ta:tangent}
A free entry $Z_{x(q),p}$ has nonzero image in the vertex cotangent space
if and only if $x(pq)\leq y$. The entries with this property generate the
actual coordinate ring $A_{x,y}$.
\end{lemma}

\begin{proof}
On the corresponding coordinate axis, each defining minor is at most
linear in the free entry. The axis satisfies the rank bounds exactly when
$x(pq)\leq y$, by the rectangle test. Otherwise one defining minor
restricts to a nonzero linear term. Since the entries have distinct root
weights, no linear combination of different entries changes this criterion.

Choose a cocharacter positive on all the roots $e_{x(q)}-e_{x(p)}$ in
question. The coordinate ring is positively graded with degree-zero part
$\C$. An entry with zero cotangent class is a sum of products of elements
of smaller positive degree. Induction on degree therefore expresses every
entry as a polynomial in the surviving entries. Taking the appropriate
torus-weight component makes the expression weight-homogeneous. This
argument takes place in the coordinate ring itself.
\end{proof}

In particular, an entry whose class in the vertex cotangent space
vanishes is zero in the coordinate ring if its root weight is not a
sum of at least two surviving root weights.
If $x(1)=y(1)$, the first-column rank bounds force that column to equal
its pivot column, and its pivot row has no other nonzero entries.
Deleting that row and column and standardizing gives an actual smaller
slice, compatibly with its Schubert strata. We may thus delete a common
prefix and assume, whenever $v<w$, that $a=v(1)<w(1)=c$.
Define
\[
 j=\min\{h>1:a<v(h)\leq c\},\qquad
 b=v(j),\quad u=v(1j),\quad t=Z_{b,1}.
\]

\begin{lemma}\label{ta:open}
The permutation $u$ satisfies $v\lessdot u\leq w$, and there is an actual
isomorphism
\[
 N_{v,w}[t^{-1}]\simeq\Gm\times N_{u,w}
\]
that preserves ordinary Schubert strata.
\end{lemma}

\begin{proof}
The choice of $j$ excludes values strictly between $a,b$ in the intervening
positions, so the exchange is a Bruhat cover. On its rectangle
$1\leq k<j$, $a<r\leq b$, the rank of $v$ counts only values greater than
$c$. The bound for $w$ contains at least those high values, by the
inequality at threshold $c+1$, as well as its first value $c$. This gives
one unit of slack and proves $u\leq w$.

More generally, for a cover $u=v(ij)$, the prefix Pl\"ucker coordinates
that define $\Omega_u$ restrict on $C^v$ to $\pm t$ for $i\leq k<j$ and
to $\pm1$ otherwise, where $t=Z_{v(j),i}$. Any other nonconstant term in
the determinant expansion would require an increasing path through an
intermediate position with value strictly between $v(i),v(j)$. The empty
rectangle of a cover excludes such a path. Consequently
$N_{v,w}\cap\Omega_u=D(t)$.

The two-sided chart isomorphism of \cite[Lemma 2.1]{KWY13} is
$\Omega_u\simeq C^u\times C_u$. Its first projection preserves the
ordinary cell of a point; its second preserves the opposite cell.
The conditions of belonging to $X_w$ and to $C^v$ therefore become,
simultaneously, conditions on the first and second factors. Restriction
of that isomorphism gives
$D(t)\simeq N_{u,w}\times(C_u\cap C^v)$.
Both sides carry their reduced scheme structures. For a cover, the second
factor is the open Richardson curve $\Gm$. This proves
the asserted isomorphism without cancelling affine factors.
\end{proof}

\subsection{The reduced first-line fibre}

Let $p_r$ be the $r$th coordinate section of the pullback of
$\OO_{\PP^{n-1}}(1)$ under the map $F_\bullet\mapsto F_1$.
In the normalized chart at $v$, $p_a=1$ and $t=p_b/p_a$.
In this cell-product chart, $p_b$ is independent of the cell factor.
Indeed, the only possible additional terms in $p_b(gM)$ have the form
$g_{bq}Z_{q,1}$, where $q>b$ and $v^{-1}(q)<j$. If $q\leq c$ this
contradicts the definition of $j$; if $q>c$, the first-line rank bound
already gives $Z_{q,1}=0$. Thus $p_b(gM)=Z_{b,1}$.

Scan the values of $w$ after its first position that are less than $c$.
Write their successive record highs as $d_1<\cdots<d_m<c$, with
$w(j_s)=d_s$ and $j_1<\cdots<j_m$. The exchanges changing the first
position in a coatom are exactly
$z_s=w(1j_s)\lessdot w$.
Their rectangles are $1\leq k<j_s$, $d_s<r\leq c$.
For $b<c$, put $H_b=\overline{C_w\cap\{p_b=0\}}\subset X_w$.

\begin{lemma}\label{ta:fibre}
The parameter divisor on $N_{v,w}$ is reduced. Its vertical components
are exactly $N_{v,z_s}$ with $v\leq z_s$ and $d_s<b$. If $b<c$, its
only other possible component is the restriction of $H_b$, which occurs
exactly when $H_b$ meets $\Omega_v$. Globally,
\[
 \begin{aligned}
  \operatorname{div}(p_b)&=H_b+\sum_{d_s<b}X_{z_s}
     &&\text{for }b<c,\\
  \operatorname{div}(p_c)&=\sum_sX_{z_s}.
 \end{aligned}
\]
The resulting divisor-class relation is
$[H_b]=\sum_{d_s\geq b}[X_{z_s}]$.
In particular, $\operatorname{div}_{N_{v,w}}(t)=\sum_{E\subset V(t)}E$,
where the sum is over its prime components, each with coefficient one.
\end{lemma}

\begin{proof}
On $C_y$, the section $p_b$ is identically zero when $y(1)<b$, nowhere
zero when $y(1)=b$, and an affine coordinate when $y(1)>b$.
The first case contributes a divisor only for a coatom, and the last
only for the big cell $C_w$. The latter zero hyperplane is irreducible.
This proves the stated component lists for the two global zero divisors.

The big-cell component has multiplicity one. At the generic point of a
vertical coatom, choose a transverse root-line parameter $\tau$ so that
the first column has the form $g(e_{d_s}+\tau e_c)$. For $d_s<b\leq c$, its $b$th
entry is $g_{bc}\tau$. If $b=c$, $g_{cc}$ is a unit. If $b<c$, the
record condition permits the independent generic cell coordinate
$g_{bc}$, so this coefficient is nonzero at the generic point. Every
vertical multiplicity is therefore one.

A nonzero section on a Cohen--Macaulay variety defines a divisor with
no embedded associated primes. Generic reducedness consequently implies
reducedness. Cell-factor independence descends the assertion to the
slice. Finally $p_b,p_c$ are sections of the same line bundle; subtracting
their divisors gives the stated relation for $[H_b]$.
\end{proof}

\subsection{Determinantal divisor classes}\label{ta:determinantal}

We use Nagata localization and polynomial invariance from
Lemma~\ref{alg:localization}, and the faithful generic-cell comparison
of Lemma~\ref{alg:generic}. The following standard calculation makes
explicit the integral local assertion needed for the detectors
\cite{Fos73,BV88}.

For positive integers $r,s$, let $D_{r,s}$ be the variety of $r\times s$
matrices of rank at most one. For $r,s\geq2$, its coordinate ring
$R=\C[x_{ij}]/I_2(x)$ is the normal Segre semigroup ring of dimension
$r+s-1$. The ruling primes $P=(x_{1j})_j$ and $Q=(x_{i1})_i$ satisfy
$\operatorname{div}(x_{11})=P+Q$. Localizing gives
\[
 R[x_{11}^{-1}]
 =\C[x_{11}^{\pm1},x_{i1}\ (i>1),x_{1j}\ (j>1)],
\]
whose units are $\C^*x_{11}^{\Z}$. The divisor localization sequence,
including its map from units, gives
\[
 \Cl(R)=(\Z[P]\oplus\Z[Q])/\Z(P+Q)\simeq\Z.
\]
For completeness, normality follows because the semigroup is the set
of pairs of nonnegative integral vectors having equal coordinate sums,
which is saturated in its group.

The same class group occurs at the vertex $\mathfrak m=(x_{ij})$.
If $n[P]$ vanishes locally for $n\neq0$, set $a=|n|$.
The homogeneous divisorial ideal $P^{(a)}$ becomes principal at
$\mathfrak m$, so $P^{(a)}/\mathfrak mP^{(a)}$ is one-dimensional.
A homogeneous lift and graded Nakayama imply that it is principal
already over $R$, contradicting the global class-group calculation.
Since localization on class groups is surjective, we obtain
\[
 \Cl(R_{\mathfrak m})\simeq\Z,\qquad [P]=1,\quad [Q]=-1.
\]
Every nonzero matrix has a nonzero-entry chart that is smooth, whereas
the vertex has tangent dimension $rs>r+s-1$. Thus
$\Sing(D_{r,s})=\{0\}$ for $r,s\geq2$.
The calculation remains valid over every field $F$ in place of $\C$:
the localized ring has units $F^*x_{11}^{\Z}$ and the graded Nakayama
argument is unchanged. In particular, this local calculation applies
to the ordinary local ring over the function field of a Schubert cell.

\subsection{Filtered maxima and actual determinantal slices}

Fix records $s<t$ and a threshold $d_t<b\leq c$.
For $y$ with $y(1)<c$, let $B(y)$ denote the first value in
$(y(1),c]$ encountered after position one. Let $\mathcal F$ be the finite
set of permutations $x\leq z_s,z_t$ such that $x\nleq z_h$ for every
$s<h<t$ and $B(x)\geq b$.

\begin{theorem}[Filtered determinantal slices]\label{ta:filtered}
Suppose $\mathcal F$ is nonempty, and let $x$ be any Bruhat-maximal
element of it. There are $p,q\geq1$ and $h\geq0$ such that
$N_{x,z_t}\simeq D_{q,p}\times\A^h$, and, as actual affine varieties,
\begin{equation}\label{ta:eq:ambient_models}
 N_{x,w}\simeq
 \begin{cases}
  D_{q+1,p},&h=0,\quad p\geq2,\\
  D_{2,p+h},&h>0,\quad q=1.
 \end{cases}
\end{equation}
In particular, $X_x$ is an irreducible component of $\Sing(X_w)$.
The first-difference parameter at $x$ becomes a matrix entry, and
$X_{z_s},X_{z_t}$ become its two primitive ruling divisors.
\end{theorem}

\begin{proof}
We supply the coordinate and rank calculations because the placement
of the two specified divisors is needed, in addition to the singularity
type.

\smallskip\noindent\emph{Normalization of the maximum.}
Write $a=d_s$ and $d=d_t$. Obtain $x_0$ from $w$ by replacing the values
$(c,a,d)$ in positions $(1,j_s,j_t)$ with $(a,d,c)$. Counting separately
before $j_s$, between $j_s,j_t$, and after $j_t$, gives
$r_{x_0}=\min(r_{z_s},r_{z_t})$. Consequently,
$y\leq z_s,z_t$ if and only if $y\leq x_0$.
For $k<j_s$, the minimum is $r_w(r,k)-\mathbf1_{a<r\leq c}$;
for $j_s\leq k<j_t$, it is $r_w(r,k)-\mathbf1_{d<r\leq c}$;
and for $k\geq j_t$ it is $r_w(r,k)$. This verifies the equality
at every threshold, including those outside the rectangles.

If $x(1)<a$, the first-difference cover for $[x,x_0]$ raises the
first value to the earliest encountered value in $(x(1),a]$.
Since $B(x)\geq b>d>a$, that value occurs after $B(x)$, so the exchange
leaves $B$ unchanged. Exclusion from intermediate coatoms is preserved
on moving upward. The resulting larger element would belong to
$\mathcal F$, contrary to maximality. Hence $x(1)=a$.

Suppose $B=B(x)<c$, at position $j$, and let $k>j$ be the first
position with $Q=x(k)\in(B,c]$. The exchange of $j,k$ is an upward
cover with rectangle $j\leq u<k$, $B<r\leq Q$. On this rectangle,
the rank of $x$ counts only values greater than $c$. Denote their
prefix counts in $x,x_0$ by $V_u,W_u$. For $u<j_t$, recordhood says
that all values of $x_0$ at most $c$ in this prefix are at most
$d<B$. At threshold $B$, the prefix of $x$ already contains $B$;
thus $x\leq x_0$ gives $V_u+1\leq W_u$. For $u\geq j_t$, the
prefix of $x_0$ contains $c$, and the comparison at threshold $c+1$
gives $V_u\leq W_u$. In either case there is one unit of slack
throughout the exchange rectangle. The exchange stays below $x_0$,
raises $B$ to $Q$, and preserves all intermediate-coatom exclusions,
again contradicting maximality. We have shown $x(1)=a$ and $B(x)=c$.
If $j=x^{-1}(c)$, this means
$x(k)\leq a\ \text{or}\ x(k)>c\;(1<k<j)$.

We also need the following consequence of recordhood. Every $y\leq w$
with $y(1)=a$ lies below $z_s$. On the coatom rectangle for $z_s$,
all values of $w$ before $j_s$ other than $c$ are less than $a$ or
greater than $c$, so $r_w(a,k)=r_w(r,k)$. Since $y(1)=a<r$,
we have
\[
 r_y(r,k)+1\leq r_y(a,k)\leq r_w(a,k)=r_w(r,k).
\]
Outside the rectangle use $y\leq w$. Maximality of $x$ consequently
excludes the simultaneous conditions
$y>x,\; y\leq z_t,\; y(1)=a,\; B(y)=c$.

\smallskip\noindent\emph{The actual support of the lower slice.}
Put $z=z_t$. A surviving tangent coordinate of $N_{x,z}$ exchanges
positions $u<v$ with $x(u)<x(v)$. If $u=1$, the new first value
cannot exceed $z(1)=d<c$; by the prefix condition its target occurs
after $j$ and lies in $(a,d]$. If $u>1$, maximality
forces the exchange to destroy $B=c$. Exchanges wholly before or after
$j$ do not do so, and neither does an exchange with $v=j$.
An exchange with $u<j<v$ must bring a value of $(a,c)$ before $j$;
hence $x(u)<a<x(v)<c$. An exchange with $u=j$ must have $x(v)>c$,
and an intermediate value in $(a,c)$ must occur between $j,v$;
otherwise the new first eligible value would still be $c$.

All surviving roots therefore have the forms $L\longrightarrow M$
and $c\longrightarrow H$, where
$L\subseteq[1,a]$, $M\subseteq[a+1,c-1]$, and
$H\subseteq[c+1,n]$. The values in $L$ occur before $j$, while those
in $M,H$ occur after $j$. The first group is nonempty and $a\in L$,
by the first-difference cover for $[x,z]$.
The negative ends of these roots belong to $L\cup\{c\}$ and their
positive ends to the disjoint set $M\cup H$. No sum of at least two
of them is a root. Lemma~\ref{ta:tangent} thus makes every other raw
entry zero in the actual ring. The matrix has only its permutation
pivots, a possibly sparse block $Z$ with rows $M$ and source columns
$L$, and entries $Y_h$ in column $j$, with rows in $H$.

\smallskip\noindent\emph{The core slack.}
For prefixes $k<j$, row $c$ is zero, so ranks at thresholds $c,c+1$
agree. Evaluating on the dense ordinary stratum of the slice gives
$j\leq j_t=z^{-1}(c)$ and
$r_z(c+1,k)=r_x(c+1,k)$ for $k<j$; high rows in these prefixes have
only fixed pivots.
If an intermediate record had $j_h\leq j$, then on its coatom
rectangle the prefix condition would give
\[
 r_x(r,k)=r_x(c+1,k)\leq r_w(c+1,k)\leq r_w(r,k)-1.
\]
The final unit is contributed by $w(1)=c$. This would imply
$x\leq z_h$, contrary to the filter. Thus $j<j_h$ for $s<h<t$.
Together with $j\leq j_t$ and the record property of $a$, this says
that $w(2),\ldots,w(j-1)$ contain no value in $(a,c]$.
In the same prefix of $z$, the sole value in $(a,c]$ is its first
value $d$. Consequently
$r_z(a+1,j-1)-r_x(a+1,j-1)=1$.
The southwest block with rows at least $a+1$ and columns at most
$j-1$ contains all of $Z$. Its source pivot rows lie in $L$ and are
absent; its target pivot columns lie after $j$ and are absent.
Every other nonzero fixed pivot has its own row and column disjoint
from $Z$. Removing these pivots from the rank condition leaves
exactly $\rank Z\leq1$.

\smallskip\noindent\emph{Recovering the whole lower slice.}
Set $p=|L|$, $q=|M|$, and $h=|H|$. Every active row and column of
$Z$ has a nonzero entry. If $Z_{ml}$ were missing, choose nonzero
entries $Z_{m'l}$ and $Z_{ml'}$. The corresponding minor would give
$Z_{m'l}Z_{ml'}=0$, contrary to integrality. Thus $Z$ is a complete
$q\times p$ matrix.

We have a closed torus-invariant embedding
$N_{x,z}\subseteq D_{q,p}\times\A^h$. The product of its surviving
coordinates is nonzero in its domain, so the slice has a point where
every $Z,Y$ entry is nonzero. All such points of the ambient product
form one torus orbit: write $Z_{ml}=u_mv_l$ and $Y_h=y_h$, with all
factors nonzero, and choose diagonal entries
\[
 t_l=v_l^{-1},\qquad t_m=u_m,\qquad t_c=1,\qquad t_h=y_h.
\]
The disjoint value sets make the choices compatible. This orbit is
dense, so a closed invariant subvariety containing one of its points
contains its closure. Therefore
$N_{x,z}\simeq D_{q,p}\times\A^h$; since $z\lessdot w$,
$\dim N_{x,w}=p+q+h$.

\smallskip\noindent\emph{A dense torus orbit in the ambient slice.}
The surviving ambient roots include the lower roots and the
first-difference root $a\longrightarrow c$. Their graph connects
$V=L\cup M\cup\{c\}\cup H$ and has root rank
$|V|-1=p+q+h$. This is the usual rank of the incidence vectors
$e_v-e_u$ of a connected graph. Choose a point $P\in N_{x,w}$ where
all surviving coordinates are nonzero. Its torus orbit has dimension
equal to the rank of the surviving weights. This dimension is at least
$p+q+h$ and at most $\dim N_{x,w}=p+q+h$; hence the orbit is dense.

Normalize each surviving coordinate $f_\alpha$ by its nonzero value
at $P$. On the orbit it then restricts to the character $\chi^\alpha$.
Because these entries generate the actual coordinate ring, we have
\[
 \C[N_{x,w}]=\C[\mathbb N\{\alpha:\alpha\text{ survives}\}].
\]
Indeed, restriction to the dense orbit is injective and its image is
exactly the algebra generated by these normalized characters.
Equal character monomials are therefore equal in the actual slice
ring.

An additional surviving root cannot involve a value outside
$V$, since that would increase root rank. Nor can a surviving root
be a sum of two others, since its normalized coordinate would then
belong to the square of the vertex ideal. The normalization multiplies
coordinates only by nonzero constants and therefore preserves their
zero ideals, including the marked divisors used below.

In the following southwest-rank calculation we retain the original
lower-slice matrix entries with their fixed pivots. The normalized
characters are used when writing the final ambient determinantal equations.

\smallskip\noindent\emph{Generic southwest pivots.}
Index the source positions by $\ell_1=1<\cdots<\ell_p<j$, the
middle positions by $j<m_1<\cdots<m_q$, and the high positions by
$j<h_1<\cdots<h_h$. Write $x(\ell_i)=\lambda_i$, $x(m_r)=\mu_r$,
and $x(h_i)=\eta_i$. Within each list the values strictly decrease:
\[
 \lambda_1=a>\cdots>\lambda_p,\qquad
 \mu_1>\cdots>\mu_q,\qquad
 \eta_1>\cdots>\eta_h>c.
\]
For instance, an increasing source pair has a reflection rectangle
contained in that from the earlier source to an active middle target.
It would give an exchange below $z$ preserving the first value and
$B=c$, contrary to maximality. The first source cannot occur
in such an increasing pair, because its value is $a$. For middle or high targets,
use an edge to the later target to contain the increasing-pair
rectangle; the exchange again preserves the same data.

For descending $\lambda_i,\mu_r$, with every $\mu_r>\lambda_i$,
consider successive source columns
$e_{\lambda_i}+v_iU$, where $U=\sum_{r=1}^q u_re_{\mu_r}$ and
$u_rv_i\neq0$, followed by target columns
$e_{\mu_1},\ldots,e_{\mu_q}$. For a prefix of $t\geq1$ source
columns, if a threshold retains $A$ source rows, its rank is
$A+\mathbf1_{\mu_1\geq r,\ A<t}$. These are the ranks of the
values $\mu_1,\lambda_1,\ldots,\lambda_{t-1}$.

After all source columns and $u<q$ target columns, if a threshold
retains $A$ source rows and $B$ target rows, its rank is
$A+\min(u,B)+\mathbf1_{A<p,\ B>u}$. These are the ranks of
$\mu_1,\ldots,\mu_{u+1},\lambda_1,\ldots,\lambda_{p-1}$.

The last target restores all original pivot rows. The generic pivot
sequence is therefore
$\mu_1,\lambda_1,\ldots,\lambda_{p-1},
\mu_2,\ldots,\mu_q,\lambda_p$.
Applying this to the dense rank-one locus of $Z$ gives
\begin{align}
 z(\ell_1)&=\mu_1,&
 z(\ell_i)&=\lambda_{i-1}\quad(2\leq i\leq p),\notag\\
 z(m_r)&=\mu_{r+1}\quad(r<q),&
 z(m_q)&=\lambda_p.\label{ta:eq:middle_pivots}
\end{align}
The star is the one-source case of the same calculation: when $h>0$,
$z(j)=\eta_1$, $z(h_i)=\eta_{i+1}$ for $i<h$, and $z(h_h)=c$.
If $h=0$, $z(j)=c$. Other positions keep their $x$-values.
Interlacing the two groups does not change the calculation because
their row and column supports are disjoint: every southwest rank is
the sum of their ranks and those of the independent fixed pivots.
The dense ordinary stratum is indexed by $z$, so these are its actual
values. In particular, $\mu_1=d$, while $j_t=j$ for $h=0$ and
$j_t=h_h$ for $h>0$. The permutation $w=z(1j_t)$ exchanges the first
value $d$ with $c$.

\smallskip\noindent\emph{The possible parameters.}
Put $K=\min(m_1,j_t)$. We claim that
\[
 r_w(r,k)-r_x(r,k)\geq1
 \qquad\text{whenever }a<r\leq c\text{ and }1\leq k<K.
\]
For $k<j$, use the high-threshold comparison and $w(1)=c$.
For $j\leq k<K$, all sources but no middle target have appeared.
Formula \eqref{ta:eq:middle_pivots} changes the ranks in $(a,c]$
by $\mathbf1_{r\leq d}$; the star does not change these counts;
and $z\to w$ contributes $\mathbf1_{d<r\leq c}$. The sum is one.
If an intermediate record occurred at a position at most $K$, its
entire coatom rectangle would be covered by
the rank bound just proved, giving $x\leq z_h$.
Thus no intermediate record occurs at or before $K$.

If $h=0$ and $p=1$, \eqref{ta:eq:middle_pivots} puts $a$ at position
$m_q>j=j_t$ in $w$, contrary to $j_s<j_t$. Hence $p\geq2$ in this
case. Now suppose $h>0$ and $q\geq2$. If $m_1<j_t$, then
$w(m_1)=\mu_2\in(a,d)$. If $a$ has not yet occurred in $w$, it
cannot subsequently be a record. Otherwise a record strictly between
$a,d$ has appeared at or before $m_1=K$, again a contradiction.
Thus $m_1>j_t$.

The surviving star edge $c\to\eta_1$ must have a value in $(a,c)$
between positions $j,h_1$; without it the exchange would preserve
$B=c$ and contradict maximality. All active middle
targets lie after $m_1>j_t\geq h_1$, so this intermediate value
belongs to neither active group and is unchanged in $w$.
It occurs before the record $d$, so it is less than $d$.
It either occurs before $a$, destroying recordhood of $a$, or forces
an intermediate record before $j_t=K$. Both are impossible; hence
$q=1$. With a unique middle position $m_1$, the same argument excludes
$m_1>h_1$: the required intermediate value would again be unchanged
and occur before $\min(m_1,j_t)=K$. We have proved that $p\geq2$
when $h=0$, while $q=1$ and $m_1<h_1$ when $h>0$.

\smallskip\noindent\emph{All ambient roots and the actual equations.}
Every $l\in L$ has an ambient edge $l\to c$. For $l=a$ this is
the first-difference edge. For $l<a$, the known rectangles of
$l\to d$ and $a\to c$ cover that of $l\to c$, because
$l<a<d<c$ and $\operatorname{pos}(l)<j<\operatorname{pos}(d)$.
There are no roots outside $V$, by the rank argument, and no
increasing roots within $L,M,H$, by the strict decrease in each value list.
No $L\to H$ root survives: it is the sum of $l\to c$ and
$c\to h$, contradicting the semigroup-algebra identity above and tangent survival.
Only $M\to H$ remains to be considered. Thus for $h=0$ the full
ambient list is $L\to M\cup\{c\}$.

For $h>0$, we have $M=\{d\}$ and its position precedes every high
target. Every $d\to\eta_i$ survives. Above threshold $c$, its
rectangle is contained in that of $c\to\eta_i$.
For $d<r\leq c$, the prefix has passed the sole middle target, so
the lower middle block has restored its original ranks, and the star
does not change these counts. The exchange $z\to w$ adds one
because $\operatorname{pos}(d)\leq k<h_i\leq j_t$.
The complete ambient list is therefore $L\to\{d,c\}\to H$.

In the case $h=0$, arrange the normalized characters in a matrix with
rows $M\cup\{c\}$ and columns $L$. Opposite products in each
$2\times2$ minor have the same character, so all such minors vanish.
Its entries generate the ring, giving a closed embedding into
$D_{q+1,p}$. Both varieties are integral of dimension $p+q$, so
this is an isomorphism. When $h>0$, arrange them as
\[
 \begin{pmatrix}
 X_{c,l}\ (l\in L)&Y_{h,d}\ (h\in H)\\
 X_{d,l}\ (l\in L)&Y_{h,c}\ (h\in H)
 \end{pmatrix}.
\]
The difference between its two row weights is always $e_c-e_d$,
so the same character argument makes its minors zero. All generators
occur and both dimensions are $p+h+1$, giving $D_{2,p+h}$.
This proves \eqref{ta:eq:ambient_models}; no affine factor has been
suppressed.

The parameter restrictions just established make both matrix dimensions at least two.
The ambient slice is singular exactly at its vertex, so the actual
cell chart gives $\Sing(X_w)\cap\Omega_x=C_x$.
The singular locus is closed and $B$-stable and therefore contains
$X_x$. A larger singular Schubert component containing $X_x$ would
have a nonempty open intersection with $\Omega_x$ of dimension
larger than $\dim C_x$, a contradiction. Thus $X_x$ is a maximal
singular component.

Finally, the prefix parameter $Z_{c,1}$ is the surviving root
$a\to c$, hence a nonzero scalar times the indicated matrix entry.
Its zero divisor has exactly two reduced ruling primes.
The distinct coatom slices $N_{x,z_s},N_{x,z_t}$ are codimension-one
subvarieties in this zero divisor because their first values are less
than $c$, so they are precisely the two rulings.
The local determinantal calculation gives their primitive opposite classes.
\end{proof}

Two marked models illustrate the cases. For $w=4231$, take the records
$a=2$, $d=3$ and threshold $b=c=4$. The filtered maximum is $x=2143$,
with $z_s=2431$, $z_t=3241$, and $(p,q,h)=(2,1,0)$.
The opposite-cell matrix and remaining equation are
\[
 \begin{pmatrix}
 0&1&0&0\\1&0&0&0\\B&D&0&1\\A&C&1&0
 \end{pmatrix},\qquad AD-BC=0.
\]
The determinantal matrix is $\bigl(\begin{smallmatrix}A&C\\B&D\end{smallmatrix}\bigr)$,
the parameter is $A$, and the two prescribed coatom ideals are
$(A,B)$ and $(A,C)$.

For the other case, take $w=3412$, $a=1$,
$d=2$, $b=c=3$. The filtered maximum is $x=1324$, with $z_s=1432$,
$z_t=2413$, and $(p,q,h)=(1,1,1)$. Its matrix is
\[
 \begin{pmatrix}
 1&0&0&0\\B&0&1&0\\A&1&0&0\\C&D&E&1
 \end{pmatrix}.
\]
The southwest rank ideal is $(C,AD+BE)$, so the actual slice has
matrix $\bigl(\begin{smallmatrix}B&A\\-D&E\end{smallmatrix}\bigr)$ of rank
at most one. Again the parameter is $A$, and the marked ideals are
$(A,B)$ and $(A,E)$. In both examples the southwest rank conditions
are exactly the displayed equations; the rank criterion also checks
the filtered maximum directly.

\subsection{Removing a horizontal component}

Return to $a=v(1)<b=v(j)\leq c=w(1)$. The vertical components of
its fibre are the active records $v\leq z_s$ with $d_s<b$.
If there are any, let $z_*=w(1h_*)$ be the one with largest first
value $d=z_*(1)<b$. We retain its failure to be Cartier while
moving to a point whose first eligible value is $c$.

\begin{lemma}[Horizontal incidence]\label{ta:incidence}
For $v\leq w$ and $b<c$, whether or not $b=B(v)$, set
$j_b=v^{-1}(b)$ and $L_b=\min\{h>1:b\leq w(h)<c\}$.
Let $V_k=\#\{h\leq k:v(h)>c\}$ and
$W_k=\#\{h\leq k:w(h)>c\}$. Then
\[
 H_b\cap\Omega_v\neq\varnothing
 \quad\Longleftrightarrow\quad
 V_k<W_k\text{ for every }j_b\leq k<L_b.
\]
An empty range imposes no condition. In particular, exchanging two
values less than $b$, while remaining below $w$, preserves incidence.
\end{lemma}

\begin{proof}
For $g\in B$, the first line of $gwB$ is generated by $ge_c$.
The inverse image of $C_w\cap H_b$ is the irreducible hypersurface
$g_{bc}=0$ in $B$. The point $gwB$ lies in $\Omega_v$ exactly when
every prefix minor $\det g_{I_k,J_k}$ is nonzero, where $I_k=v([k])$
and $J_k=w([k])$.
The intersection $\Omega_v\cap H_b$ is nonempty exactly when it meets the
dense subset $H_b\cap C_w$. On an irreducible hypersurface, finitely
many minors are simultaneously nonzero somewhere exactly when none
vanishes identically there.

For $I=I_k$, $J=J_k$, determinant monomials correspond to perfect
matchings in the bipartite graph with row vertices $I$, column vertices
$J$, and edges $r\leq s$. The entries, including invertible diagonal
entries, are independent, so distinct matching monomials do not cancel.
There is a matching because $v\leq w$. Deleting the edge $(b,c)$
can violate a Hall inequality only for a row subset $S$ containing
$b$ that loses column $c$ from its neighborhood. Every other row of
$S$ must exceed $c$. Its neighborhood after deletion is then
$(J\cap[b,n])\setminus\{c\}$, independent of those other rows.
The largest possible $S$ is $\{b\}\cup(I\cap(c,n])$.
Since the original Hall inequalities hold, failure occurs precisely
when $b\in I$ and $1+\#(I>c)=\#(J\cap[b,n])$.
Because $c\in J$ and $\#(I>c)\leq\#(J>c)$, this is equivalent to
\[
 b\in I,\qquad J\cap[b,c)=\varnothing,\qquad
 \#(I>c)=\#(J>c).
\]
These conditions say $k\geq j_b$, $k<L_b$, and $V_k=W_k$.
Their absence is precisely the claimed incidence criterion.
A swap of two values less than $b$ changes neither $j_b$ nor the
high-value counts.
\end{proof}

\begin{lemma}\label{ta:two_components}
If the first-difference fibre has both a horizontal and a vertical
component, there is $x$ with $v\leq x\leq z_*$, $x(1)=d$, and
$B(x)=b$ whose fibre has precisely the components $z_*$ and $H_b$.
The divisor $z_*$ is not Cartier at the slice vertex.
\end{lemma}

\begin{proof}
While the first value $a_0$ is less than $d$, exchange it with the
earliest value in $(a_0,d]$. The first-difference cover for the upper
bound $z_*$ keeps the new permutation below $z_*$. Since $b>d$ is
the first value in $(a_0,c]$, the exchanged value occurs after $b$.
The exchange preserves the first eligible value $b$ and involves two
values less than $b$; Lemma~\ref{ta:incidence} preserves horizontal
incidence. Lemma~\ref{ta:fibre} applies at each step.
The first value strictly increases until it reaches $d$.
Lower records are now excluded by the first-column Bruhat bound;
records in $(d,b)$ were initially inactive and cannot become active
on moving upward. The sole remaining vertical component is $z_*$.

In a normal Noetherian local domain, if
$t\in\mathfrak m\setminus\mathfrak m^2$ and
$\operatorname{div}(t)=E_1+\cdots+E_r$ is reduced, no nonempty proper
subsum is Cartier. Otherwise that effective subsum is
$\operatorname{div}(f)$ with $f\in\mathfrak m$ by normality; its
effective nonempty complement is $\operatorname{div}(t/f)$ with
$t/f\in\mathfrak m$, forcing $t\in\mathfrak m^2$.
Here the parameter survives in the tangent space by
Lemma~\ref{ta:tangent}. Both components are invariant for the
contracting torus and contain the vertex. Applying the observation to
$\operatorname{div}(t_b)=z_*+H_b$ proves the last assertion.
\end{proof}

\begin{lemma}[The reduced first-column prime]\label{ta:column}
Let $d$ be a first-position record of $w$, let $z_*=w(1h_*)$ with
$w(h_*)=d$, and suppose $y\leq w$ and $y(1)=d$.
Then $y\leq z_*$. In $N_{y,w}$, the reduced prime ideal of
$N_{y,z_*}$ is $J_d=(Z_{p,1}:p>d)$. At the vertex its minimal
number of generators is $\mu(J_d)=q(y,w)$, where
\[
 q(y,w):=\#\{k>1:y(k)>d,\ y(1k)\leq w\}.
\]
It is Cartier there if and only if $q(y,w)=1$.
\end{lemma}

\begin{proof}
Relative to $w$, $z_*$ lowers the ranks by one on
$1\leq k<h_*$, $d<r\leq c$. If $H_k$ counts values greater than
$c$ in the corresponding prefix of $w$, recordhood gives
$r_w(d,k)=H_k+1$. Since $y(1)=d<r$,
$r_y(r,k)+1\leq r_y(d,k)\leq H_k+1$, proving $y\leq z_*$.

Consider the Schubert variety $Y_d=\{F_1\subset E_d\}$.
The support of $X_w\cap Y_d$ is the union of first-position record
coatoms with first value at most $d$. To see this, take $x\leq w$
with $a_0=x(1)\leq d$ and choose the first record $d_0\geq a_0$,
at position $h$. Such a record exists and $d_0\leq d$.
Before $h$, all nonfirst values of $w$ are less than $a_0$ or greater
than $c$. On the coatom rectangle of $w(1h)$, the inequality
$r_x(r,k)+1\leq r_x(a_0,k)\leq r_w(a_0,k)$ supplies the unit of
slack needed for $x\leq w(1h)$. The reverse support inclusion is
immediate.

We also need a scheme-theoretic assertion. Intersections of Schubert
varieties are reduced; the compatible Frobenius-splitting result and
its characteristic-zero form are recalled in
\cite[Section 2, p.~2474]{Kum17}. On $\Omega_y$, lower records are
absent because $y(1)=d$. Thus the intersection is scheme-theoretically
$X_{z_*}\cap\Omega_y$. The condition $F_1\subset E_d$ has ideal
generated by first-column entries in rows $p>d$.
Multiplication by the upper triangular cell factor acts on these tail
entries by an invertible triangular matrix and preserves their ideal.
The reduced intersection descends through the cell-product chart to the
asserted equality $J_d=(Z_{p,1}:p>d)$ of reduced prime ideals.
Its height is one.

By Lemma~\ref{ta:tangent}, the surviving first-column entries are
exactly those counted by $q(y,w)$. Their distinct weights make their
cotangent classes linearly independent. They generate all of $J_d$:
express any first-column entry of weight $e_p-e_d$ as a
weight-homogeneous polynomial in the surviving entries.
View a monomial as a nonnegative integral flow along edges from
smaller to larger values. A flow of weight $e_p-e_d$ is a path from
$d$ to $p$: inflow equals outflow elsewhere and there are no directed
cycles. It contains an edge with source $d$, which, since $y(1)=d$,
is a surviving first-column entry. Every monomial belongs to their
ideal. Their classes remain independent in $J_d/\mathfrak mJ_d$
because $\mathfrak mJ_d\subseteq\mathfrak m^2$, proving
the asserted generator count $\mu(J_d)=q(y,w)$. A height-one prime in a normal local
domain is Cartier exactly when its divisorial ideal is principal.
\end{proof}

\begin{lemma}[Greedy prefix lift]\label{ta:greedy}
Suppose $y\leq w$, $y(1)=d<c$, and $q(y,w)\geq2$.
There is $y'\geq y$ with $y'\leq w$, $y'(1)=d$, $B(y')=c$, and
$q(y',w)\geq2$. If $d$ is a record with associated coatom $z_*$,
all permutations in the construction lie below $z_*$.
\end{lemma}

\begin{proof}
If $B(y)=c$ there is nothing to prove. Otherwise put
\[
 \begin{aligned}
  j&=\min\{h>1:d<y(h)\leq c\},&
  B&=y(j)<c,\\
  k&=\min\{h>j:B<y(h)\leq c\},&
  Q&=y(k),\qquad \widetilde y=y(jk).
 \end{aligned}
\]
The position $k$ exists because $c$ occurs
after $j$. No intervening value lies strictly between $B,Q$, so this
is an upward cover with rectangle
$j\leq s<k,\; B<r\leq Q$.
On this rectangle every value of $y$ at most $c$ in the prefix is
less than $r$, so its rank counts only values above $c$.
The threshold-$c+1$ comparison with $w$, together with the first
value $c$, supplies one unit of slack. Therefore
$\widetilde y\leq w$, its first value is $d$, and its first eligible
value is $Q>B$.

Call $T>d$ a target if exchanging its position with position one
stays below $w$. Every target is at most $c$ by the first-line bound;
$B$ is a target by Lemma~\ref{ta:open}. Choose a second target
$T\neq B$. The new first-difference target is $Q$; we exhibit
another.

If $T<B$, put $h=y^{-1}(T)$ and $Y=y(1h)\leq w$.
Here $h>j$ and $h\neq k$, so $\widetilde y(1h)=Y(jk)$.
On this rectangle, all values at most $c$ in $Y$
remain below $r$, including the new first value $T$ and the value
$d$ at $h$. The high-threshold comparison for $Y\leq w$ and the
first value $c$ of $w$ again supply one unit of slack.
Thus $Y(jk)\leq w$, and $T<Q$ remains a target.

If $T>B$, put $h=y^{-1}(T)$; the definition of $k$ gives $h\geq k$.
Set $Y_T=y(1h)\leq w$, $Y_B=y(1j)\leq w$, and
$X=\widetilde y(1k)$. At positions $(1,j,k)$, the values of
$y,\widetilde y,X$ are respectively $(d,B,Q)$, $(d,Q,B)$, and
$(B,Q,d)$. For a prefix $s<j$, $X$ has the same values as $Y_B$.
For $j\leq s<k$ and $d<r\leq B$, its rank is the rank of $y$
plus one, equal to that of $Y_T$ because $h\geq k$ and $T>B$.
At all other thresholds in this middle prefix range, its rank equals
that of $\widetilde y$. For $s\geq k$, $X$ and $\widetilde y$
have the same prefix value sets. Every rank of $X$ is bounded by
that of one of $Y_B,Y_T,\widetilde y$, all below $w$.
Hence $X\leq w$, so $B<Q$ is the other target.

Each lift increases $B$, preserves $d$ and $q\geq2$, and remains
below $w$. Iteration terminates at $B=c$. When $d$ is a record,
Lemma~\ref{ta:column} puts every step below $z_*$.
\end{proof}

\begin{proposition}[Horizontal-to-prefix witness]\label{ta:horizontal}
Suppose the original first-difference fibre has both a horizontal and
a vertical component. There is $x\geq v$ with $x(1)=d$, $B(x)=c$,
and $x\leq z_*$ such that its prefix fibre has the following
properties: its sole active record of first value less than $b$ is
$z_*$; it has an active record of first value at least $b$; and
$z_*$ is not Cartier at the slice vertex.
\end{proposition}

\begin{proof}
Lemma~\ref{ta:two_components} gives a point of first value $d$ with
two-component fibre $z_*+H_b$ and with $z_*$ non-Cartier.
Lemma~\ref{ta:column} gives $q\geq2$, and Lemma~\ref{ta:greedy}
raises $B$ to $c$ while retaining this inequality and the coatom
upper bound. Thus $z_*$ remains non-Cartier at the endpoint.

All moves are upward, so inactive records cannot become active.
The first value $d$ excludes lower records, leaving $z_*$ as the
only record below $b$. The prefix fibre has no horizontal component,
by Lemma~\ref{ta:fibre}. If $z_*$ were its only component,
reducedness would give $\operatorname{div}(p_c/p_d)=z_*$, making
it Cartier, a contradiction. A further active record therefore exists,
and its first value is at least $b$.
\end{proof}

\subsection{Integral coefficient detectors and induction}

Fix $V(t)$ in $N_{v,w}$. Order its active vertical components
$E_1,\ldots,E_k$ by record first value, and place the horizontal
component $H$, if present, last. A class supported on this fibre has
a representative
\[
 D=\sum_{i=1}^k c_i[E_i]+c_H[H],\qquad c_i,c_H\in\Z,
\]
omitting the horizontal term when absent. This representation need
not be unique.

\begin{lemma}[Coefficient differences]\label{ta:differences}
For each adjacent pair $E_i,E_{i+1}$, there is a maximal singular
Schubert component $X_y$ containing $C_v$ at whose generic point the
local image of $D$ is, up to choice of a
primitive generator, $\pm(c_i-c_{i+1})\in\Z$.
If $H$ and a vertical component both occur, there is such a component
detecting $\pm(c_k-c_H)$.
\end{lemma}

\begin{proof}
Let $z_s,z_t$ correspond to adjacent verticals. At $v$, no intervening
record is active and $B(v)=b>d_t$, so $v$ belongs to
the filtered set $\mathcal F$ of Theorem~\ref{ta:filtered}. Choose a maximal element $y$ of that
set above $v$ and apply Theorem~\ref{ta:filtered}.
Its prefix fibre is the zero divisor of one matrix entry, with
exactly the rulings $z_s,z_t$. Lemma~\ref{ta:fibre} then excludes
every other active record in that chart. If the original $H_b$
exists, its class is the sum of high records $d_h\geq b$ by
the class relation in Lemma~\ref{ta:fibre}; all are absent from this chart.
The ordinary detector ring is
\[
 \OO_{X_w,\eta_{X_y}}
 \simeq(\C(C_y)[z_{ij}]/I_2(z))_{(z_{ij})},
\]
for the matrix size supplied by Theorem~\ref{ta:filtered}.
The two ruling classes are primitive opposites by the calculation in
Section~\ref{ta:determinantal}, which applies over the field $\C(C_y)$.
This gives the stated integer difference. If $P$ is the corresponding
slice prime in $A_{v,w}$ and $R=(A_{v,w})_P$, the cell-product chart also
identifies this detector with $R(T_1,\ldots,T_{\ell(v)})$.
The comparison $R\to R(T_1,\ldots,T_{\ell(v)})$ is faithfully flat and local,
so Lemma~\ref{alg:generic} reflects vanishing.

For the last difference, Proposition~\ref{ta:horizontal} gives
$x\geq v$ with first value $d$ and $B(x)=c$, where only the highest
original low record $z_*=E_k$ remains and at least one high record
is active. Choose the active high record $z_t$ of smallest first
value. It is adjacent to $z_*$ in this prefix fibre.
Use the filtered construction with threshold $d_t+1\leq c$ and a
maximal element $y\geq x$. Its two rulings are exactly $z_*,z_t$.
Every other original low record is absent, while
the class relation in Lemma~\ref{ta:fibre} restricts to
$[H_b]=[X_{z_t}]=-[X_{z_*}]$. The local coefficient of the original
class is consequently $\pm(c_k-c_H)$.
This divisor-class relation supplies the comparison without requiring
an equation for the old horizontal divisor in the new chart.
\end{proof}

\begin{theorem}[Affine slice detection]\label{ta:affine}
For every $v\leq w\in S_n$, restriction induces an injection
\begin{equation}\label{ta:eq:affine_injection}
 \Cl(N_{v,w})\longrightarrow
 \bigoplus_{S\in\Irr\Sing(N_{v,w})}
 \Cl(\OO_{N_{v,w},\eta_S}).
\end{equation}
\end{theorem}

\begin{proof}
Induct on $\ell(w)-\ell(v)$; the zero-dimensional case is immediate.
Delete a common prefix and choose the first-difference cover and
parameter from Lemma~\ref{ta:open}.
Let $D$ map to zero on the right of \eqref{ta:eq:affine_injection}.
A nonempty open subset retains exactly the irreducible singular
components meeting it; it does not turn a deeper singular stratum
into a new component. By Lemma~\ref{ta:open} and Laurent
polynomial invariance, the class on $D(t)$ comes from
$D'\in\Cl(N_{u,w})$. The maximal singular components of this
product are $\Gm\times S'$. At their generic points,
Lemma~\ref{alg:generic} shows that $D'$ vanishes at all maximal
singular generic points of $N_{u,w}$. Since its dimension is one less,
induction gives $D'=0$.

Lemma~\ref{alg:localization} writes $D$ as an integer combination of
the vertical classes $[E_i]$ and, when present, the horizontal class
$[H]$. Every detector in
Lemma~\ref{ta:differences} is a singular component of the original
slice: its global maximal singular Schubert component meets
$\Omega_v$, and its intersection is the product of $C_v$ with
its slice factor. The vanishing hypothesis therefore gives
$c_i=c_{i+1}$ for $i<k$, and $c_k=c_H$ when both kinds occur.
All coefficients equal some $m\in\Z$.
This is automatic for a single fibre component, including a sole
horizontal component. Reducedness from Lemma~\ref{ta:fibre} now gives
\[
 D=m\sum_{E\subset V(t)}[E]=m[\operatorname{div}(t)]=0.\qedhere
\]
\end{proof}

\begin{proof}[Proof of Theorem~\ref{ta:main}]
Take a scheme point $p\in C_v\subset X_w$ and
$\alpha\in\Cl(\OO_{X_w,p})$ vanishing at all maximal singular
generic points whose closures contain $p$. Put
\[
 A_{\mathrm{ch}}=\C[\Omega_v\cap X_w]
 \simeq A_{v,w}[t_1,\ldots,t_{\ell(v)}].
\]
Localization on class groups is surjective by taking closures of the
height-one primes in a representative. Choose a lift
$\beta\in\Cl(A_{\mathrm{ch}})$ of $\alpha$.

Every maximal singular component is a Schubert variety $X_z$.
It meets $\Omega_v$ exactly when $v\leq z$; in that case
the cell-product chart shows that it contains all of $C_v$, hence $p$
even if $p$ is not closed. Its generic point is a generalization
of $p$, so the restriction of $\beta$ factors through $\alpha$.
Thus every chosen lift vanishes at all maximal singular generic points
of this chart.

Polynomial invariance writes $\beta$ as the pullback of
$\gamma\in\Cl(A_{v,w})$. The full-cell generic local ring is
$R(T)$ for the corresponding slice generic local ring $R$.
Lemma~\ref{alg:generic} therefore makes $\gamma$ vanish at every
maximal singular generic point of the slice.
Theorem~\ref{ta:affine} gives $\gamma=0$, hence $\beta=0$ and
$\alpha=0$. This proves \eqref{ta:eq:local_injection} without any
equivariance assumption on $\alpha$.

The invertible locus of a finite rank-one reflexive sheaf is open.
Non-Cartierness at $\eta_S$ therefore persists on $S$, proving one
inclusion in \eqref{ta:eq:locus}. If $p$ avoids the indicated union,
the local class of $D$ vanishes at every relevant generic point;
the injection makes it zero at $p$, proving the other inclusion.

Finally, for $P\supset B$, the projection $\pi:G/B\to G/P$ is
locally trivial with irreducible smooth fibre $P/B$.
The inverse image of a partial-flag Schubert variety is the full-flag
Schubert variety indexed by the longest representative of its coset.
Its singular locus is the inverse image of the singular locus
downstairs, and inverse images of irreducible components are
irreducible. Weil divisors pull back on these local products;
Cartierness is equivalent before and after the faithfully flat smooth
extension. The generic-point maps are compatible with $\pi$.
Apply \eqref{ta:eq:locus} upstairs and descend to obtain the
partial-flag assertion.
\end{proof}
